\subsection{赋范空间中的绝对可和族}

定义 8. 在赋范空间 E 中，E 的点的族 \((\pmb{x}_{t})\) 称为绝对可和的，如果范数的族 \((||\pmb{x}_{t}||)\) （诸 \(\pmb{x}_{t}\) 的范数所成的族）在 R 中可和。

这个概念看起来依赖于 E 上范数的选取；但由第 3 号的命题 7 与实数可和族的比较原理，关于 E 上范数 p 绝对可和的族，关于 E 上任何与 p 等价的范数也绝对可和。

若 \((\pmb{x}_t)_{i\in I}\) 是 \(E\) 中点的既可和又绝对可和的族，则我们有

\[
\left| \left| \sum_ {i \in I} x _ {i} \right| \right| \leqslant \sum_ {i \in I} | | x _ {i} | |.\tag{4}
\]

事实上，对 I 的每个有限子集 J 有 \(\left\|\sum_{i\in J}x_{i}\right\| \leqslant \sum_{i\in J}||x_{i}||\) ，对 I 的有限子集组成的有向集取极限即得不等式 (4)。

命题 11. 在完备赋范空间 E 中，每个绝对可和族都可和。

因为若 \((\pmb{x}_i)\) 是 \(\mathbf{E}\) 中的绝对可和族，则对每个 \(\varepsilon > 0\) ，有一个有限子集 \(J\) （指标集 \(I\) 的），使得对每个有限子集 \(H\) （\(I\) 中的、不与 \(J\) 相交的），有 \(\sum_{i \in H} ||\pmb{x}_i|| \leqslant \varepsilon\) ；更有 \(\left\| \sum_{i \in H} \pmb{x}_i \right\| \leqslant \varepsilon\) ，由于 \(E\) 完备，这就证明了命题（柯西准则，第 III 章，§ 5，第 2 号，定理 1）。

在 E 中，通项为 \(x_{n}\) 的级数称为绝对收敛的，如果通项为 \(\|x_{n}\|\) 的级数在 R 中收敛，或（等价地）如果族 \((x_{n})\) 绝对可和；于是（第 III 章，§ 5，第 7 号，命题 9）：

推论. 在完备赋范空间 E 中，每个绝对收敛级数都交换收敛。

命题 11 的逆命题一般不成立。

例如考虑空间 \(\mathcal{B}(\mathbf{N};\mathbf{R})\) （实数的有界序列 \(x = (x_{n})_{n\in \mathbb{N}}\) 所成，带范数 \(\| x\| = \sup_n |x_n|\) ）。设 \(x_{m}\) 是序列 \((x_{mn})_n\in \mathbb{N}\) ，它满足：\(x_{mn} = 0\) （若 \(m\neq n\) ），且 \(x_{mm} = 1 / m\) （对 \(m\geqslant 1\) ）。立即可验证，序列 \((x_m)_{m\in \mathbb{N}}\) 在 \(\mathcal{B}(\mathbf{N};\mathbf{R})\) 中可和，其和是元素 \(y = (y_{n})\) ，满足 \(y_0 = 0\) 且 \(y_{n} = 1 / n\) （若 \(n\geqslant 1\) ）；但由于 \(\| x_m\| = 1 / m\) ，诸 \(x_{m}\) 的范数组成的序列在 \(\mathbf{R}\) 中不可和。

然而，我们在第 VII 章，§ 3，第 1 号中已看到，\(\mathbf{R}^n\) 中每个可和族都绝对可和。

